\textbf{5.\ \textsc{Ferromagnetic stripe} (12 points)} --- \textit{Jaan Kalda, Eero Uustalu.}

\textbf{Tools} A caliper, a ruler, graph papers, a resistive magnetic field
sensor connected to batteries in a battery holder, a multimeter with two wires, a
magnet a stripe made from soft ferromagnetic material of thickness
$0.25\,\mathrm{mm}$ --- \textbf{do not bend excessively to avoid damaging}.

\textbf{i)} \textit{(0.5 points)}Connect the banana ends of the wires to the COM
port and to the V$\Omega$mA-port of the multimeter. Switch on the multimeter in
20 volt (DC) range, and touch the two metallic leads of the battery holder (which
are next to the points were the red and black wires come out from the holder)
with the crocodile ends of the wires. Record the voltage $\mathcal{E}$ on the
output leads of the battery holder. If the voltage is below $3.0\,\mathrm{V}$,
you may ask for replacement batteries.

For all your magnetic field measurements, keep in mind that \textit{if the
battery voltage were to be exactly $3\,\mathrm{V}$, each millivolt in the reading
would correspond to 10 microteslas of the magnetic field strength. However, the
reading in millivolts is proportional to both the magnetic field and to the
battery voltage}.

Connect the crocodiles to the yellow and red wires of the magnetic sensor. Keep
in mind that (a) the sensor may have a non-zero offset: even if there is no
magnetic field, the multimeter reading $V_0$ might be non-zero; (b) there is
always the magnetic field of Earth. In what follows, avoid measuring magnetic
fields which correspond to voltage readings bigger than $500\,\mathrm{mV}$ ---
such strong fields may cause changes in the offset value $V_0$. If you
accidentally expose sensor to such fields, determine and use the new value of
$V_0$.

The magnetic sensor has a small white dot marked on one of its edges. This points
to the direction of that magnetic field component which is being measured.

\textbf{ii)} \textit{(1.5 points)}Determine the offset voltage $V_0$ and the
magnitude of the Earth's magnetic field $B_E \equiv |\vec{B}_E|$, and the angle
between the vertical direction and the direction of $\vec{B}_E$.

Now, attach the magnet to the ferromagnetic stripe so that its circular face is
touching the stripe's surface near one of its ends. Let us use a perpendicular
system of coordinates where $x - y$-plane is the plane of the stripe, with the
longest symmetry axis of the stripe serving as the $x$-axis, and $x = 0$ being
at the position of the centre of the magnet.

The total magnetic field is the superposition of the field of the permanent
magnet $\vec{B}_m$, the field of the magnetised ferromagnetic stripe $\vec{B}$,
and the Earth's magnetic field $\vec{B}_E$. Below we are interested only in
$\vec{B}$. Assume that $\vec{B}_m$ depends only on the distance from the magnet
and remains unchanged when the magnet is detached from the stripe.

\textbf{iii)} \textit{(2.5 points)}Measure the vertical field
$B_z = B_z(L/2{,}y)$ caused by the stripe with the magnet, as a function of $y$,
for $-w/2 \leq y \leq w/2$, at $x = L/2$, where $w$ denotes the width and $L$
--- the length of the stripe. Find the ratio $\kappa = \langle B_z\rangle$ to
$B_z(L/2{,}0)$, where the average magnetic field
\[
\langle B_z\rangle \equiv \int_{-w/2}^{w/2} B_z(L/2{,}y)\mathrm{d}y.
\]

Assume that $\kappa$ remains constant along the stripe.

\textbf{iv)} \textit{(3.5 points)} Measure $B_z(x{,}0)$ near the surface of the
stripe, as a function of $x$, and plot the measurement results.

\textbf{v)} \textit{(2.5 points)} Let $J_s$ denote the saturation magnetisation
of the stripe material; estimate the value of $J_s\mu_0$ (this is, roughly
speaking, the strongest magnetic B-field which the ferromagnet is able to
carry).

\textbf{vi)} \textit{(1.5 points)} Prove experimentally that by small values of
$x$, the magnetisation inside the stripe has reached saturation.
